% 連続体力学 連続体力学-1 問題3（003-kijutsu） 解いた日 2026-10-19
%% kotae
\toi{1}{◎}{}
質点…大きさがない。

剛体…相対位置が変わらない。

連続体…相対位置が変わる。

歪み…相対位置の変化量
\toi{2}{△}{立式}
固定…E

流れ…L

移流項　$\ayamari{(\vect{v}\cdot\vect{\nabla})\vect{v}}{(\vect{u}\cdot\vect{\nabla})T}$
\naze{$(\vect{v}\cdot\vect{\nabla})\vect{v}$ は速度場に物質微分を当てたとき（加速度）の移流項である。温度計の読みの差は式(4)で $A=T$ とした $\mathrm{D}T/\mathrm{D}t-\partial T/\partial t=(\vect{u}\cdot\vect{\nabla})T=\sum_{i}u_{i}\,\partial T/\partial x_{i}$ である。}
\toi{3}{◎}{}
ある。各点で温度が変化しないだけで、流れとともに温度計が違う温度の点に行ったら変化する。
%% kizuki
\item \mae{第1回 問題3・10-17}移流項は「何の量の変化率か」で中身が変わる。温度なら $(\vect{u}\cdot\vect{\nabla})T$，加速度なら $(\vect{u}\cdot\vect{\nabla})\vect{u}$ である。前回は演算子 $\vect{u}\cdot\vect{\nabla}$ だけを書き，今回は速度に当てた形を書いた。式(4)の $A$ に問われている量（ここでは $T$）を入れてから書く。
\item 速度の記号は教材に合わせて $\vect{u}$ にする（$\vect{v}$ と混ぜない）。
